\documentclass[a4paper,12pt]{article}

\usepackage[T2A]{fontenc}
\usepackage[utf8]{inputenc}
\usepackage[russian]{babel}
\usepackage{amsmath,amssymb,amsfonts}
\usepackage{graphicx}
\usepackage{float}
\usepackage{caption}
\usepackage{subcaption}
\usepackage{geometry}
\usepackage{setspace}
\usepackage{indentfirst}
\usepackage{enumitem}
\usepackage{booktabs}
\usepackage{algorithm}
\usepackage{algpseudocode}
\usepackage{hyperref}
\usepackage{amsthm}
\usepackage{thmtools}
\usepackage{xcolor}

\geometry{left=3cm, right=1.5cm, top=2cm, bottom=2cm}
\onehalfspacing
\setlength{\parindent}{1.25cm}
\setlength{\parskip}{0pt}
\hypersetup{
    colorlinks=true,
    linkcolor=blue,
    filecolor=magenta,
    citecolor=red,
    urlcolor=cyan,
    pdftitle={Курсовой проект},
    pdfpagemode=FullScreen,
}

\declaretheoremstyle[
    spaceabove=6pt, spacebelow=6pt,
    headfont=\normalfont\bfseries,
    notefont=\mdseries, notebraces={(}{)},
    bodyfont=\normalfont,
    postheadspace=1em,
    headpunct={}
]{mystyle}

\declaretheorem[style=mystyle, name=Определение, numberwithin=section]{definition}
\declaretheorem[style=mystyle, name=Теорема, sibling=definition]{theorem}
\declaretheorem[style=mystyle, name=Лемма, sibling=definition]{lemma}
\declaretheorem[style=mystyle, name=Утверждение, sibling=definition]{proposition}
\declaretheorem[style=mystyle, name=Пример, numbered=no]{example}
\declaretheorem[style=mystyle, name=Замечание, numbered=no]{remark}
\renewcommand{\theHdefinition}{\thesection.\arabic{definition}}

\floatname{algorithm}{Алгоритм}
\algrenewcommand\algorithmicrequire{\textbf{Вход:}}
\algrenewcommand\algorithmicensure{\textbf{Выход:}}

\newcommand{\R}{\mathbb{R}}
\newcommand{\norm}[1]{\left\lVert #1 \right\rVert}
\newcommand{\fro}[1]{\left\lVert #1 \right\rVert_{\mathrm F}}
\newcommand{\nuc}[1]{\left\lVert #1 \right\rVert_*}
\newcommand{\tr}{\operatorname{tr}}
\newcommand{\rank}{\operatorname{rank}}
\newcommand{\grad}{\operatorname{grad}}
\renewcommand{\Proj}{\mathcal{P}}
\newcommand{\Mcal}{\mathcal{M}}
\newcommand{\Retr}{\operatorname{Retr}}
\newcommand{\St}{\mathrm{St}}
\newcommand{\Gr}{\mathrm{Gr}}

\begin{document}

\begin{titlepage}
    \centering
    {\large\bfseries Национальный исследовательский университет\\
    «Высшая школа экономики»\par}
    \vspace{1cm}
    {\large Факультет экономических наук, Факультет компьютерных наук\par}
    {\large Образовательные программы «Экономика и анализ данных», «Прикладная математика и информатика»\par}
    \vspace{2cm}
    {\Large\bfseries КУРСОВОЙ ПРОЕКТ\par}
    \vspace{1cm}
    {\Large\bfseries Геометрические методы на многообразии матриц фиксированного ранга для заполнения матриц\par}
    \vspace{2cm}
    \begin{flushright}
        \large
        \textbf{Выполнили:}\\
        Студент 2 курса, группы БЭАД245\\
        Аляутдинов Карим Маратович,\\
        Студент 2 курса, группы БПМИИ245\\
        Исмаилов Магомед Айдынович.\\
        \vspace{0.5cm}
        \textbf{Научный руководитель:}\\
        профессор, д.ф.-м.н.\\
        Широков Дмитрий Сергеевич\\
    \end{flushright}
    \vfill
    {\large Москва, 2026\par}
\end{titlepage}

\tableofcontents
\newpage

\section*{Введение}
\phantomsection
\addcontentsline{toc}{section}{Введение}


\textbf{Актуальность темы.} Задача восстановления пропусков в матрицах возникает во многих прикладных задачах: в рекомендательных системах, обработке изображений, восстановлении панельных экономических и социальных данных. Часто именно восстановленная таблица становится основой дальнейшего анализа, поэтому важно понимать не только итоговую ошибку метода, но и его устойчивость к шуму, количеству пропусков, структуре наблюдений и выбору ранга.

В экономических данных такая задача выглядит особенно естественно. Показатели стран, регионов, компаний или временных периодов часто зависят от небольшого числа скрытых факторов. Это приводит к гипотезе низкого ранга: большая матрица может приближенно описываться через несколько латентных компонент. На этой идее строится задача \textit{matrix completion}, то есть восстановление недостающих элементов по наблюдаемой части матрицы.

Классические подходы к matrix completion часто используют выпуклую релаксацию через ядерную норму. У таких методов есть сильная теория, но в больших задачах они могут быть тяжелыми вычислительно, поскольку требуют сингулярных разложений и других дорогих операций. Факторизационные методы вроде ALS обычно работают быстрее, но сильнее зависят от инициализации и параметров. На этом фоне риманов градиентный спуск интересен как промежуточный путь: он работает напрямую на множестве матриц фиксированного ранга и использует геометрию этого множества при построении шага. Особенно важна компактная реализация через $X=USV^\top$, поскольку она позволяет не делать полное SVD большой матрицы на каждой итерации.

В этой работе мы рассматриваем тему не только как реализацию одного алгоритма. Основной интерес состоит в том, чтобы понять поведение риманова градиентного спуска в разных режимах: при шуме, разной доле наблюдаемых элементов, разных инициализациях и неверном выборе ранга. Такой подход позволяет сделать из курсовой не просто обзор метода, а небольшой вычислительный исследовательский проект.

\textbf{Цель исследования.} Реализовать и протестировать метод риманова градиентного спуска на многообразии матриц фиксированного ранга для задачи заполнения матриц, а затем сравнить его качество, скорость и устойчивость с базовыми методами matrix completion.

\textbf{Задачи исследования.} Для достижения цели в работе решаются следующие задачи:
\begin{enumerate}[label=\arabic*.]
    \item провести обзор литературы по выпуклым, факторизационным и геометрическим методам заполнения матриц;
    \item описать геометрическую постановку задачи на многообразии матриц фиксированного ранга;
    \item вывести формулу риманова градиента и шаг риманова градиентного спуска для функции невязки по наблюдаемым элементам;
    \item рассмотреть регуляризованную версию риманова спуска для шумных данных;
    \item реализовать экспериментальный код для риманова градиентного спуска, его регуляризованной и компактной $USV^\top$-версии, ALS и Soft-Impute как базовых методов сравнения;
    \item подготовить серию синтетических экспериментов с контролируемым рангом, шумом и долей пропусков;
    \item исследовать устойчивость методов к шуму, доле наблюдаемых элементов, выбору ранга и инициализации;
    \item проанализировать скорость сходимости и вычислительную эффективность алгоритмов;
    \item на следующем этапе добавить прикладную проверку на экономических или панельных данных и расширить сравнение с OptSpace.
\end{enumerate}

\textbf{Объект и предмет исследования.}
\begin{itemize}
    \item \textbf{Объект исследования:} процессы восстановления пропущенных значений в низкоранговых матрицах.
    \item \textbf{Предмет исследования:} риманов градиентный спуск на многообразии матриц фиксированного ранга и его свойства в сравнении с альтернативными алгоритмами matrix completion.
\end{itemize}

\textbf{Методы исследования.} В работе используются методы линейной алгебры, элементы дифференциальной геометрии и римановой оптимизации, анализ литературы, а также вычислительный эксперимент. Качество восстановления предполагается оценивать по RMSE, MAE и относительной ошибке в норме Фробениуса, а вычислительную эффективность --- по времени работы, числу итераций и траекториям сходимости.

\textbf{Структура работы.} В первом разделе представлен обзор литературы по matrix completion. Во втором разделе вводятся основные обозначения. В третьем разделе приводится теоретическая часть: сначала кратко вводятся необходимые геометрические понятия, затем формулируется задача на многообразии матриц фиксированного ранга и описывается риманов градиентный спуск. В четвертом разделе описываются реализация и дизайн вычислительных экспериментов. В заключении формулируются промежуточные выводы и дальнейшие шаги работы.

\newpage

\section{Обзор литературы}
\label{sec:literature}

\subsection{Классическая выпуклая линия}

Современная теория matrix completion начинается с выпуклой постановки через ядерную норму. В работе Candes и Recht \cite{candes-recht} показано, что матрицу малого ранга можно восстанавливать по случайно наблюдаемым элементам, если заменить минимизацию ранга на выпуклую релаксацию. Для литературы это был принципиальный шаг: задача перестала быть чисто эвристической и получила строгую математическую основу.

Дальше эта линия была существенно усилена. В работе Candes и Tao \cite{candes-tao} получены почти оптимальные по порядку оценки на число наблюдений, достаточное для восстановления. Статья Recht \cite{recht} дала более компактное и прозрачное доказательство близких результатов. Если смотреть на эти три работы вместе, то они задают базовую выпуклую логику всей темы: сначала показывается, что восстановление вообще возможно, затем уточняются условия, при которых это можно гарантировать.

Для нашей работы эта линия важна по двум причинам. Во-первых, она дает естественный baseline: если мы хотим оценивать новый или более сложный метод, то сравнивать его разумно именно с выпуклой постановкой, у которой есть сильная теория. Во-вторых, эти статьи показывают ограничения такого подхода. Основной акцент в них делается на гарантии восстановления, тогда как вычислительная цена решения в больших задачах остается отдельным вопросом.

\subsection{Алгоритмы для больших матриц}

После теоретических результатов почти сразу возникает вычислительная проблема. Даже если выпуклая модель выглядит естественно, ее надо уметь считать на больших матрицах. В работе Cai, Candes и Shen \cite{cai-svt} предложен алгоритм Singular Value Thresholding, который переводит общую выпуклую постановку в последовательность более управляемых шагов.

Эту же линию продолжает работа Mazumder, Hastie и Tibshirani \cite{mazumder}, где предлагается Soft-Impute. Этот алгоритм особенно важен для прикладной части темы, потому что ориентирован на большие и разреженные матрицы. В отличие от более абстрактного теоретического языка первых статей, здесь уже явно виден переход к реальным вычислительным задачам.

Для нашей работы эти алгоритмы нужны не просто как исторические примеры. Скорее, это естественные конкуренты геометрического подхода. Если риманов градиентный спуск должен быть содержательно интересен, он должен сравниваться не с абстрактной выпуклой моделью, а с реально используемыми численными методами.

\subsection{Шум и более реалистичная постановка}

Классические результаты о точном восстановлении важны, но в прикладной задаче этого недостаточно. В реальных данных шум почти неизбежен: наблюдаемые значения сами по себе измеряются с ошибкой, а шаблон пропусков часто далек от идеального случайного. Работа Candes и Plan \cite{candes-plan} важна именно тем, что переводит matrix completion в постановку с шумом.

Для нашего проекта это ключевой момент. Если ограничиться только безошибочной моделью, то значительная часть будущего сравнения алгоритмов просто теряет смысл. Нас интересует как раз то, что происходит при реальных возмущениях: устойчивость к шуму, качество восстановления скрытых наблюдений и чувствительность метода к меняющимся условиям эксперимента.

\subsection{Невыпуклые и фиксированно-ранговые методы}

Хотя выпуклая линия очень сильна теоретически, она не всегда оптимальна вычислительно. Это естественным образом приводит к невыпуклым постановкам, где ранг задается заранее, а матрица ищется в факторизованной форме. Одной из ранних и важных работ в этом направлении является статья Keshavan, Montanari и Oh \cite{keshavan}, где предложен метод OptSpace.

Для нас OptSpace важен не только как конкретный алгоритм, но и как переходный этап в развитии темы. Здесь исследователь уже не заменяет ранг выпуклой оболочкой, а работает непосредственно с низкоранговой структурой. Именно такая логика и подводит к геометрическим методам на многообразиях фиксированного ранга.

\subsection{Геометрический подход}

Прямой геометрический подход к matrix completion подробно развивается в работе Vandereycken \cite{vandereycken}. В этой статье задача формулируется как оптимизация функции невязки на многообразии матриц фиксированного ранга, а все основные объекты --- касательное пространство, риманов градиент, ретракция --- вводятся в естественной для многообразия форме. Для нашей темы эта статья является одной из центральных, потому что именно она задает базовую схему риманова градиентного метода, который мы планируем реализовать.

Работа Mishra и Sepulchre \cite{r3mc} развивает ту же линию, но уже с более сложной геометрической конструкцией. Авторы предлагают алгоритм R3MC и подчеркивают, что выбор подходящей римановой метрики может заметно повлиять на устойчивость и численное поведение метода. Для нашего проекта это особенно полезно, потому что мы заранее хотим рассматривать не только итоговую ошибку восстановления, но и более тонкие характеристики: зависимость от старта, чувствительность к шуму и скорость сходимости.

Важное отличие геометрической линии от выпуклой состоит в том, что здесь акцент смещается с вопроса ``когда восстановление гарантировано'' на вопрос ``как именно ведет себя алгоритм при прямой оптимизации на низкоранговом множестве''. Именно поэтому геометрические методы хорошо подходят для проекта с исследовательским уклоном.

\subsection{Расширения и прикладной контекст}

Литература по matrix completion давно вышла за рамки одной базовой постановки. В работе Jain и Dhillon \cite{inductive} рассматривается \textit{inductive matrix completion}, где кроме самой матрицы используются и дополнительные признаки строк и столбцов. Это важно, потому что в прикладных данных часто есть внешняя информация, которую можно использовать при восстановлении.

Статья Athey et al. \cite{causal-panel} показывает, что методы matrix completion уже применяются и в задачах с панельными экономическими данными. Для нашей работы это служит прикладным ориентиром: тема не ограничивается абстрактной линейной алгеброй, а реально используется в задачах, близких к экономике и анализу панелей.

Наконец, обзор Davenport и Romberg \cite{overview} полезен тем, что он связывает разные направления low-rank recovery в одну картину. Из него хорошо видно, что выбор метода почти всегда связан с компромиссом между теоретической строгостью, вычислительной стоимостью и гибкостью модели.

\subsection{Выводы из обзора литературы}

Из обзора литературы для нашей работы следуют несколько важных выводов. Во-первых, выпуклые методы с ядерной нормой остаются естественным baseline, потому что они хорошо изучены и понятны теоретически. Во-вторых, методы фиксированного ранга и геометрические алгоритмы особенно интересны в больших задачах, где прямая работа с низкоранговой структурой может дать вычислительное преимущество. В-третьих, именно для геометрических методов особенно важны практические вопросы, которые не сводятся к одной только гарантии восстановления: устойчивость к шуму, зависимость от инициализации, чувствительность к выбору ранга и поведение по ходу итераций.

Именно здесь и находится место нашей работы. Мы не пытаемся доказывать новые теоремы о восстановлении, но и не ограничиваемся простой реализацией известного алгоритма. Основная идея проекта состоит в том, чтобы реализовать риманов градиентный спуск на многообразии матриц фиксированного ранга и систематически сравнить его с альтернативами по нескольким критериям. Такая постановка позволяет превратить курсовую работу в небольшой исследовательский проект.

\newpage

\section{Основные определения и обозначения}
\label{sec:definitions}

Пусть дана матрица
\[
Y \in \R^{m \times n},
\]
в которой наблюдаются только элементы с индексами из множества
\[
\Omega \subseteq \{1,\dots,m\} \times \{1,\dots,n\}.
\]
Предполагается, что существует скрытая матрица
\[
X_* \in \R^{m \times n},
\]
имеющая малый ранг, и наблюдения удовлетворяют соотношению
\[
Y_{ij} = (X_*)_{ij} + \varepsilon_{ij}, \qquad (i,j)\in\Omega,
\]
где $\varepsilon_{ij}$ описывает шум.

\begin{definition}
Введем маску наблюдений $W\in\R^{m\times n}$:
\[
W_{ij} =
\begin{cases}
1, & (i,j)\in\Omega,\\
0, & (i,j)\notin\Omega.
\end{cases}
\]
Тогда оператор наблюдаемых элементов удобно записывать через адамарово произведение:
\[
\Proj_\Omega(X)=W\odot X.
\]
Здесь $\odot$ означает покомпонентное произведение матриц.
\end{definition}

Тогда задача matrix completion состоит в восстановлении матрицы $X_*$ по наблюдаемой части $\Proj_\Omega(Y)$ при предположении, что искомая матрица имеет малый ранг.

\begin{definition}
Если $X = U\Sigma V^\top$ --- сингулярное разложение матрицы ранга $r$, то условие \textit{инкогерентности} означает, что сингулярные подпространства матрицы не слишком сильно сосредоточены на отдельных координатах. Именно это условие фигурирует в классических теоремах о восстановлении \cite{candes-recht,candes-tao,recht}.
\end{definition}

В дальнейшем используются следующие обозначения:
\begin{itemize}
    \item $\fro{X}$ --- норма Фробениуса матрицы $X$;
    \item $\nuc{X}$ --- ядерная норма матрицы $X$;
    \item $\rank(X)$ --- ранг матрицы $X$;
    \item $\Mcal_r = \{X \in \R^{m\times n} : \rank(X)=r\}$ --- множество матриц фиксированного ранга $r$.
\end{itemize}

\newpage


\section{Теоретическая часть}
\label{sec:theory}

\subsection{От локальных координат к многообразиям}

В этом разделе вводится минимальный геометрический аппарат, необходимый для постановки задачи на многообразии матриц фиксированного ранга. Материал основан на лекционных и семинарских материалах по дифференциальной геометрии, в частности на курсе А. О. Иванова с мехмата МГУ, а также на курсах Фоменко и Пенского \cite{ivanov-dg,fomenko-dg,penskoy-seminars}. Мы используем только те определения и факты, которые дальше нужны для риманова градиентного спуска.

\begin{definition}
Пусть $M$ --- топологическое пространство. \textit{Картой} размерности $d$ называется пара $(U,\varphi)$, где $U\subset M$ --- открытое множество, а
\[
\varphi:U\to \varphi(U)\subset \R^d
\]
--- гомеоморфизм на открытое множество $\varphi(U)$.
\end{definition}

\begin{definition}
Набор карт $\mathcal A=\{(U_\alpha,\varphi_\alpha)\}$ называется \textit{атласом}, если множества $U_\alpha$ покрывают $M$:
\[
M=\bigcup_\alpha U_\alpha.
\]
Если все карты имеют значения в $\R^d$, то $d$ называется размерностью многообразия.
\end{definition}

\begin{definition}
Топологическое пространство $M$ называется \textit{$d$-мерным топологическим многообразием}, если у каждой точки $p\in M$ есть карта $(U,\varphi)$ размерности $d$, то есть окрестность $U$ точки $p$ гомеоморфна открытому множеству в $\R^d$.
\end{definition}

Когда две карты пересекаются, одна и та же точка получает две системы координат. Поэтому нужно понимать, как координаты переходят друг в друга. Если $U_\alpha\cap U_\beta\ne\varnothing$, то отображение перехода имеет вид
\[
\varphi_\beta\circ\varphi_\alpha^{-1}:
\varphi_\alpha(U_\alpha\cap U_\beta)
\to
\varphi_\beta(U_\alpha\cap U_\beta).
\]

\begin{definition}
Атлас называется \textit{гладким}, если все отображения перехода между картами гладкие. Пространство $M$ вместе с таким атласом называется \textit{гладким многообразием}.
\end{definition}

\begin{definition}
Пусть $M$ и $N$ --- гладкие многообразия, а $F:M\to N$ --- отображение. Оно называется \textit{гладким}, если в любых картах $(U,\varphi)$ на $M$ и $(V,\psi)$ на $N$ координатное представление
\[
\psi\circ F\circ\varphi^{-1}
\]
является гладким отображением между открытыми подмножествами евклидовых пространств.
\end{definition}

Если $F$ взаимно однозначно, гладко и обратное отображение $F^{-1}$ тоже гладко, то $F$ называется \textit{диффеоморфизмом}. В этом случае многообразия считаются одинаковыми с точки зрения гладкой геометрии.

Композиция гладких отображений снова гладкая; это обеспечивает согласованность при переходе между картами.

\subsection{Вложенные многообразия и уравнения}

В задачах оптимизации допустимое множество часто лежит в $\R^N$ и задается ограничениями. Поэтому нам прежде всего нужны вложенные многообразия.

\begin{definition}
Подмножество $M\subset\R^N$ называется \textit{вложенным гладким многообразием} размерности $d$, если для каждой точки $x\in M$ существует окрестность $W\subset\R^N$ и гладкая замена координат $\Phi:W\to\Phi(W)\subset\R^N$ с гладким обратным отображением такая, что
\[
\Phi(W\cap M)=\Phi(W)\cap\bigl(\R^d\times\{0\}^{N-d}\bigr).
\]
\end{definition}

Перед теоремой о множествах, заданных уравнениями, введем регулярные и особые точки.

\begin{definition}
Пусть $F:\R^N\to\R^k$ --- гладкое отображение. Точка $x\in\R^N$ называется \textit{регулярной точкой} отображения $F$, если дифференциал $DF(x)$ имеет максимальный ранг:
\[
\rank DF(x)=k.
\]
Если это условие не выполнено, то $x$ называется \textit{особой точкой}.
\end{definition}

\begin{definition}
Точка $y\in\R^k$ называется \textit{регулярным значением} отображения $F$, если каждая точка $x\in F^{-1}(y)$ является регулярной. Если в прообразе $F^{-1}(y)$ есть особая точка, то $y$ называется \textit{особым значением}.
\end{definition}

\begin{theorem}[Теорема о регулярном уровне]
\label{thm:regular-level}
Пусть $F:\R^N\to\R^k$ --- гладкое отображение, а $0\in\R^k$ является регулярным значением $F$. Тогда множество
\[
M=F^{-1}(0)
\]
является гладким вложенным многообразием размерности $N-k$.
\end{theorem}

\begin{proof}
Это прямое следствие теоремы о неявной функции. Условие регулярности означает, что $k$ уравнений локально независимы, поэтому $k$ координат можно выразить как гладкие функции от оставшихся $N-k$ координат. Эти свободные координаты и дают локальную карту на $M$.
\end{proof}

Обратно, вложенное многообразие локально можно задать независимой системой уравнений.

\begin{theorem}[Локальное задание вложенного многообразия уравнениями]
\label{thm:local-equations}
Пусть $M\subset\R^N$ --- вложенное гладкое многообразие размерности $d$. Тогда для каждой точки $x\in M$ существует окрестность $W\subset\R^N$ и гладкое отображение
\[
F:W\to\R^{N-d}
\]
такое, что
\[
M\cap W=F^{-1}(0),
\qquad
\rank DF(z)=N-d
\quad\text{для всех } z\in M\cap W.
\]
\end{theorem}

Доказательство этого факта мы не приводим: для нас важен вывод, что вложенные многообразия можно локально понимать как множества решений независимых гладких уравнений.

В компактной реализации будут встречаться матрицы с ортонормированными столбцами:
\[
\St(r,m)=\{U\in\R^{m\times r}:U^\top U=I_r\}.
\]
Это многообразие Stiefel. Многообразие Grassmann $\Gr(r,m)$ описывает сами $r$-мерные подпространства; в нашей задаче оно появляется через столбцовое и строковое подпространства матрицы фиксированного ранга.

\subsection{Касательные векторы и касательные пространства}

В выбранной карте $(U,\varphi)$, содержащей точку $p\in M$, касательный вектор можно сначала задать как обычный вектор
\[
v_\varphi\in\R^d.
\]
Если $(V,\psi)$ --- другая карта около $p$, координаты того же вектора пересчитываются по правилу
\begin{equation}
\label{eq:tangent-coordinate-change}
v_\psi
=
D(\psi\circ\varphi^{-1})(\varphi(p))\,v_\varphi.
\end{equation}

\begin{definition}
Касательный вектор к гладкому многообразию $M$ в точке $p$ --- это согласованный набор координатных представлений $\{v_\varphi\}$ по всем картам, содержащим $p$, причем при переходе от одной карты к другой координаты связаны формулой \eqref{eq:tangent-coordinate-change}.
\end{definition}

Множество всех касательных векторов в точке $p$ образует линейное пространство $T_pM$.

Эквивалентно, касательный вектор можно понимать как скорость гладкой кривой. Если $\gamma:(-\varepsilon,\varepsilon)\to M$ и $\gamma(0)=p$, то
\[
\left.\frac{d}{dt}\right|_{t=0}\varphi(\gamma(t))
\]
задает его координаты в карте $\varphi$.

Если $M\subset\R^N$ --- вложенное многообразие, то это описание принимает привычный вид:
\[
T_xM=
\{\dot\gamma(0):\gamma(t)\subset M,\ \gamma(0)=x\}.
\]

\begin{proposition}
\label{prop:tangent-regular-level}
Если $M=F^{-1}(0)$ --- регулярный уровень из теоремы \ref{thm:regular-level}, то
\[
T_xM=\ker DF(x).
\]
\end{proposition}

\begin{proof}
Пусть $\gamma(t)\subset M$ и $\gamma(0)=x$. Тогда $F(\gamma(t))=0$ для всех малых $t$. Дифференцируя по $t$ при $t=0$, получаем
\[
DF(x)[\dot\gamma(0)]=0.
\]
Значит, $T_xM\subset\ker DF(x)$. Обратное включение следует из локальной параметризации, полученной из теоремы о неявной функции: каждый вектор из ядра дифференциала является производной некоторой кривой на уровне $F=0$.
\end{proof}

Дальше эта схема используется для множества матриц фиксированного ранга: сначала описывается допустимое множество, затем его касательное пространство, а затем градиентный шаг.

\subsection{Риманова метрика, дифференциал и градиент}

\begin{definition}
\textit{Риманова метрика} на гладком многообразии $M$ --- это выбор скалярного произведения
\[
g_p:T_pM\times T_pM\to\R
\]
в каждой точке $p\in M$, причем для любых $\xi,\eta,\zeta\in T_pM$ и $a,b\in\R$ выполнены:
\[
g_p(a\xi+b\eta,\zeta)=ag_p(\xi,\zeta)+bg_p(\eta,\zeta),
\]
\[
g_p(\xi,\eta)=g_p(\eta,\xi),
\]
\[
g_p(\xi,\xi)>0\quad\text{при}\quad \xi\ne 0.
\]
Кроме того, метрика должна гладко зависеть от точки $p$.
\end{definition}

В этой работе используется метрика, индуцированная евклидовым скалярным произведением на пространстве матриц:
\[
g_X(\xi,\eta)=\langle \xi,\eta\rangle
=\tr(\xi^\top\eta).
\]

\begin{definition}
Пусть $f:M\to\R$ --- гладкая функция. Дифференциал функции $f$ в точке $p$ --- это линейное отображение
\[
df_p:T_pM\to\R,
\]
которое на касательном векторе $\xi=\dot\gamma(0)$ задается как
\[
df_p[\xi]=\left.\frac{d}{dt}\right|_{t=0}f(\gamma(t)).
\]
\end{definition}

\begin{definition}
Риманов градиент функции $f$ в точке $p$ --- это единственный вектор $\grad f(p)\in T_pM$, такой что
\[
g_p(\grad f(p),\xi)=df_p[\xi]
\qquad
\text{для всех } \xi\in T_pM.
\]
\end{definition}

\begin{proposition}
\label{prop:projected-gradient}
Пусть $M\subset\R^N$ --- вложенное многообразие с индуцированной евклидовой метрикой, а $\bar f$ --- гладкое продолжение функции $f$ на окрестность $M$ в $\R^N$. Тогда
\begin{equation}
\label{eq:riemannian-gradient-projection}
\grad f(x)=\Proj_{T_xM}\bigl(\nabla \bar f(x)\bigr).
\end{equation}
\end{proposition}

Нормальная компонента евклидова градиента не двигает точку вдоль многообразия, поэтому в римановом методе остается только касательная часть.

\subsection{Геодезические и ретракции}

В евклидовом пространстве градиентный шаг имеет вид
\[
x_{k+1}=x_k-\eta_k\nabla f(x_k).
\]
На многообразии касательный шаг нужно вернуть обратно на допустимое множество. Через геодезическую это записывается как
\[
x_{k+1}=\operatorname{Exp}_{x_k}\bigl(-\eta_k\grad f(x_k)\bigr).
\]
На практике вместо экспоненты используют ретракции.

\begin{definition}
\label{def:retraction}
Ретракцией называется отображение
\[
\Retr_p:T_pM\to M,
\]
которое удовлетворяет условиям
\[
\Retr_p(0_p)=p,
\qquad
D\Retr_p(0_p)=\operatorname{id}_{T_pM}.
\]
\end{definition}

Ретракция совпадает с геодезическим шагом в первом порядке малости, но обычно считается проще.

Общая схема риманова градиентного спуска имеет вид
\begin{equation}
\label{eq:rgd-general}
x_{k+1}=\Retr_{x_k}\bigl(-\eta_k\grad f(x_k)\bigr).
\end{equation}
Длина шага $\eta_k$ может быть фиксированной или подбираться процедурой backtracking line search.

\subsection{Переход к задаче заполнения матриц}

Пусть дана матрица наблюдений $Y\in\R^{m\times n}$, но известны только элементы с индексами из множества
\[
\Omega\subseteq \{1,\ldots,m\}\times\{1,\ldots,n\}.
\]
Введем маску наблюдений $W\in\R^{m\times n}$:
\[
W_{ij}=
\begin{cases}
1, & (i,j)\in\Omega,\\
0, & (i,j)\notin\Omega.
\end{cases}
\]
Тогда оператор наблюдаемых элементов записывается как
\[
\Proj_\Omega(X)=W\odot X,
\]
где $\odot$ --- адамарово, то есть покомпонентное, произведение.

Классическая выпуклая постановка:
\begin{equation}
\label{eq:convex}
\min_{X \in \R^{m\times n}}
\frac12 \fro{\Proj_\Omega(X-Y)}^2 + \lambda \nuc{X}.
\end{equation}

Геометрическая fixed-rank постановка:
\begin{equation}
\label{eq:fixed-rank}
\min_{X \in \R^{m\times n}} \frac12 \fro{\Proj_\Omega(X-Y)}^2
\qquad \text{при условии} \qquad \rank(X)=r.
\end{equation}
Оптимизация ведется по множеству
\[
\Mcal_r=\{X\in\R^{m\times n}:\rank(X)=r\}.
\]

Для шумных данных используем регуляризованную версию:
\begin{equation}
\label{eq:fixed-rank-regularized}
\min_{X \in \Mcal_r} f_\lambda(X)
= \frac12 \fro{\Proj_\Omega(X-Y)}^2 + \frac{\lambda}{2}\fro{X}^2,
\qquad \lambda\ge 0.
\end{equation}
Штраф $\fro{X}^2$ стабилизирует метод при шуме и завышенном выбранном ранге.

\subsection{Многообразие матриц фиксированного ранга}

\begin{proposition}
\label{prop:fixed-rank-manifold}
Множество
\[
\Mcal_r=\{X\in\R^{m\times n}:\rank(X)=r\}
\]
является гладким вложенным многообразием размерности
\[
\dim \Mcal_r=r(m+n-r).
\]
\end{proposition}

\begin{proof}
Рассмотрим матрицу $X\in\Mcal_r$. У нее есть невырожденный минор размера $r\times r$. Перестановкой строк и столбцов можно считать, что этот минор находится в левом верхнем углу. Тогда в некоторой окрестности $X$ любую матрицу ранга $r$ можно записать блочно:
\[
\begin{pmatrix}
A & B\\
C & D
\end{pmatrix},
\]
где $A\in\R^{r\times r}$ невырождена. Условие ранга $r$ эквивалентно равенству
\[
D=CA^{-1}B.
\]
Значит, локально свободными параметрами являются блоки $A$, $B$ и $C$. Их общее число равно
\[
r^2+r(n-r)+(m-r)r=r(m+n-r).
\]
Так как зависимый блок $D$ выражается через свободные параметры гладко, получаем локальные координаты на $\Mcal_r$.
\end{proof}

Важно различать
\[
\{X:\rank(X)=r\}
\qquad\text{и}\qquad
\{X:\rank(X)\le r\}.
\]
Первое является гладким многообразием. Второе содержит матрицы меньших рангов и в целом имеет особенности, поэтому риманов метод формулируется на страте фиксированного ранга.

Если $X\in\Mcal_r$, то удобно использовать компактное сингулярное разложение
\[
X=U\Sigma V^\top,
\]
где $U\in\R^{m\times r}$ и $V\in\R^{n\times r}$ имеют ортонормированные столбцы, а $\Sigma\in\R^{r\times r}$ невырождена.

\begin{proposition}
\label{prop:fixed-rank-tangent}
Если $X=U\Sigma V^\top\in\Mcal_r$, то касательное пространство к $\Mcal_r$ в точке $X$ имеет вид
\[
T_X\Mcal_r =
\left\{
UMV^\top + U_pV^\top + UV_p^\top :
M \in \R^{r\times r},\,
U^\top U_p = 0,\,
V^\top V_p = 0
\right\}.
\]
\end{proposition}

Три слагаемых отвечают за изменение центрального блока, столбцового подпространства и строкового подпространства.

\begin{lemma}
\label{lemma:tangent-projection}
Пусть $P_U=UU^\top$ и $P_V=VV^\top$. Тогда ортогональная проекция произвольной матрицы $Z\in\R^{m\times n}$ на $T_X\Mcal_r$ имеет вид
\begin{equation}
\label{eq:tangent-projection-fixed-rank}
\Proj_{T_X\Mcal_r}(Z)=P_UZ+ZP_V-P_UZP_V.
\end{equation}
\end{lemma}

Для ее получения разложим $Z$ по столбцовому пространству $U$ и строковому пространству $V$:
\[
Z=P_UZP_V+P_UZ(I-P_V)+(I-P_U)ZP_V+(I-P_U)Z(I-P_V).
\]
Последнее слагаемое является нормальной компонентой. Удаляя его, получаем
\[
Z-(I-P_U)Z(I-P_V)=P_UZ+ZP_V-P_UZP_V.
\]

\subsection{Градиент функции потерь на многообразии фиксированного ранга}

Рассмотрим функцию
\[
f(X)=\frac12\fro{\Proj_\Omega(X-Y)}^2.
\]

\begin{lemma}
\label{lemma:euclidean-gradient}
Евклидов градиент функции $f$ равен
\[
\nabla f(X)=\Proj_\Omega(X-Y).
\]
\end{lemma}

Действительно, для любого приращения $H\in\R^{m\times n}$ линейная часть изменения функции равна
\[
\langle \Proj_\Omega(X-Y),\Proj_\Omega(H)\rangle.
\]
Так как $\Proj_\Omega$ является ортогональной проекцией,
\[
\langle \Proj_\Omega(X-Y),\Proj_\Omega(H)\rangle
=
\langle \Proj_\Omega(X-Y),H\rangle.
\]
Отсюда следует формула для евклидова градиента.

Для регуляризованной функции \eqref{eq:fixed-rank-regularized} имеем
\[
\nabla f_\lambda(X)=\Proj_\Omega(X-Y)+\lambda X.
\]
По предложению \ref{prop:projected-gradient} риманов градиент получается проекцией этого выражения на касательное пространство:
\[
\grad_{\Mcal_r} f_\lambda(X)
=
\Proj_{T_X\Mcal_r}\bigl(\Proj_\Omega(X-Y)+\lambda X\bigr).
\]

\subsection{Ретракция для матриц фиксированного ранга}

После касательного шага нужно вернуться на $\Mcal_r$. Используем ретракцию через лучшую аппроксимацию ранга $r$:
\begin{equation}
\label{eq:fixed-rank-retraction}
R_X(\Xi)=\Pi_r(X+\Xi),
\end{equation}
где $\Pi_r$ обозначает усеченное SVD, то есть ближайшую матрицу ранга $r$ в норме Фробениуса.

\begin{proposition}
\label{prop:svd-retraction}
Отображение $R_X(\Xi)=\Pi_r(X+\Xi)$ является ретракцией на $\Mcal_r$ в окрестности $X$.
\end{proposition}

Условие $R_X(0)=X$ выполнено, так как $X$ уже имеет ранг $r$. Для касательного направления $\Xi$ матрица $X+t\Xi$ отличается от некоторой кривой на $\Mcal_r$ только на $O(t^2)$. Усеченное SVD выбирает ближайшую матрицу ранга $r$, поэтому
\[
\Pi_r(X+t\Xi)=X+t\Xi+O(t^2),
\]
то есть $DR_X(0)[\Xi]=\Xi$.

Геодезики на $\Mcal_r$ вычислять сложнее, а для градиентного метода достаточно возвращения на многообразие с правильным первым порядком. Поэтому используем ретракцию.

Итак, прямой риманов градиентный спуск имеет вид
\begin{equation}
\label{eq:rgd}
X_{k+1}
=
R_{X_k}\bigl(-\eta_k\grad_{\Mcal_r} f_\lambda(X_k)\bigr).
\end{equation}

\subsection{Компактная реализация через факторы}

Ретракция \eqref{eq:fixed-rank-retraction} требует SVD матрицы размера $m\times n$ на каждой итерации. Поэтому храним текущую точку в виде
\[
X=USV^\top,\qquad U^\top U=I_r,\quad V^\top V=I_r,
\]
где $U\in\R^{m\times r}$, $V\in\R^{n\times r}$, а $S\in\R^{r\times r}$.

Это не обычный градиентный спуск по независимым переменным $U,S,V$: шаг по-прежнему является римановым шагом по матрице $X$.

Пусть
\[
G=\nabla f_\lambda(X)=\Proj_\Omega(X-Y)+\lambda X.
\]
Тогда компоненты проекции на касательное пространство можно записать так:
\[
M=U^\top G V,
\]
\[
U_\perp=GV-UM,\qquad V_\perp=G^\top U-VM^\top.
\]
В этих обозначениях
\[
\grad_{\Mcal_r} f_\lambda(X)
=UMV^\top+U_\perp V^\top+UV_\perp^\top.
\]
Матрица $M$ меняет центральный блок, а $U_\perp$ и $V_\perp$ --- столбцовое и строковое подпространства.

Пусть длина шага равна $\eta$. Обозначим
\[
\Delta M=-\eta M,\qquad
\Delta U=-\eta U_\perp,\qquad
\Delta V=-\eta V_\perp.
\]
Тогда касательный шаг записывается как
\begin{equation}
\label{eq:compact-tangent-step}
X+\Xi
=U(S+\Delta M)V^\top+\Delta U V^\top+U\Delta V^\top.
\end{equation}

Сделаем QR-разложения
\[
\Delta U=Q_UR_U,\qquad \Delta V=Q_VR_V.
\]
Тогда
\[
X+\Xi
=
\begin{bmatrix}U & Q_U\end{bmatrix}
\begin{bmatrix}
S+\Delta M & R_V^\top\\
R_U & 0
\end{bmatrix}
\begin{bmatrix}V & Q_V\end{bmatrix}^\top.
\]
Матрица $X+\Xi$ представлена через малое ядро
\[
K=
\begin{bmatrix}
S+\Delta M & R_V^\top\\
R_U & 0
\end{bmatrix},
\]
размером не больше $2r\times 2r$. Для ретракции достаточно SVD этого ядра:
\[
K=\widetilde U\widetilde S\widetilde V^\top.
\]
После усечения до ранга $r$ получаем новые факторы:
\[
U_+=\begin{bmatrix}U & Q_U\end{bmatrix}\widetilde U_r,\qquad
S_+=\widetilde S_r,\qquad
V_+=\begin{bmatrix}V & Q_V\end{bmatrix}\widetilde V_r.
\]
Итоговая новая точка равна
\[
X_+=U_+S_+V_+^\top.
\]

\begin{algorithm}[H]
\caption{Компактный риманов градиентный спуск на $\Mcal_r$}
\label{alg:compact-rgd}
\begin{algorithmic}[1]
\Require наблюдаемая матрица $Y$, маска $\Omega$, ранг $r$, параметр $\lambda$, начальные факторы $U_0,S_0,V_0$
\Ensure приближение $\widehat X=USV^\top$ ранга $r$
\For{$k=0,1,2,\ldots$}
    \State $X_k \gets U_kS_kV_k^\top$
    \State $G_k \gets \Proj_\Omega(X_k-Y)+\lambda X_k$
    \State $M_k \gets U_k^\top G_kV_k$
    \State $U_{\perp,k} \gets G_kV_k-U_kM_k$
    \State $V_{\perp,k} \gets G_k^\top U_k-V_kM_k^\top$
    \State выбрать $\eta_k$ с помощью backtracking line search
    \State $\Delta M_k\gets-\eta_kM_k$, $\Delta U_k\gets-\eta_kU_{\perp,k}$, $\Delta V_k\gets-\eta_kV_{\perp,k}$
    \State выполнить QR-разложения $\Delta U_k=Q_UR_U$, $\Delta V_k=Q_VR_V$
    \State $K_k\gets
    \begin{bmatrix}
    S_k+\Delta M_k & R_V^\top\\
    R_U & 0
    \end{bmatrix}$
    \State вычислить SVD малого ядра $K_k=\widetilde U\widetilde S\widetilde V^\top$
    \State $U_{k+1}\gets [U_k,Q_U]\widetilde U_r$
    \State $S_{k+1}\gets \widetilde S_r$
    \State $V_{k+1}\gets [V_k,Q_V]\widetilde V_r$
    \If{выполнен критерий остановки}
        \State \textbf{break}
    \EndIf
\EndFor
\State \Return $\widehat X=U_kS_kV_k^\top$
\end{algorithmic}
\end{algorithm}

Компактная реализация сохраняет геометрию fixed-rank многообразия, но заменяет SVD большой матрицы на QR-разложения тонких матриц и SVD малого ядра.


\iffalse

\subsection{Геометрические предварительные сведения}

В этой работе дифференциальная геометрия используется не как отдельная теория, а как язык для описания ограниченной оптимизации. В обычном градиентном спуске переменная может двигаться по всему линейному пространству. В нашей задаче это не так: матрица должна сохранять фиксированный ранг. Поэтому допустимое множество уже не является линейным пространством, и обычный евклидов шаг может вывести нас за пределы нужного класса матриц.

\begin{definition}
Пусть $M\subset\R^N$ --- гладкое вложенное многообразие и $x\in M$. Касательное пространство в точке $x$ определяется как
\[
T_xM=\{\dot\gamma(0):\ \gamma:(-\varepsilon,\varepsilon)\to M,\ \gamma(0)=x\},
\]
где $\gamma$ пробегает все гладкие кривые на $M$, проходящие через $x$.
\end{definition}

Это определение хорошо согласуется с интуицией из дифференциальной геометрии: касательное пространство является линейной аппроксимацией многообразия в малой окрестности точки. В задачах оптимизации оно описывает все допустимые мгновенные направления движения.

\begin{proposition}[Регулярный уровень]
Пусть $F:\R^N\to\R^k$ --- гладкое отображение, причем $\rank DF_x=k$ для всех $x\in F^{-1}(0)$. Тогда множество
\[
M=F^{-1}(0)
\]
является гладким вложенным многообразием размерности $N-k$, а его касательное пространство имеет вид
\[
T_xM=\ker DF_x.
\]
\end{proposition}

\begin{proof}
Идея доказательства следует из теоремы о неявной функции. Условие $\rank DF_x=k$ означает, что $k$ уравнений локально независимы. Поэтому в окрестности каждой точки из $M$ можно выразить $k$ координат через остальные $N-k$ координат. Отсюда $M$ локально параметризуется $N-k$ переменными. Если $\gamma(t)\subset M$ и $\gamma(0)=x$, то $F(\gamma(t))=0$. Дифференцируя по $t$ при $t=0$, получаем $DF_x\dot\gamma(0)=0$, то есть $T_xM\subset \ker DF_x$. Обратное включение следует из той же локальной параметризации: каждый вектор из ядра дифференциала реализуется как скорость некоторой кривой на уровне $F=0$.
\end{proof}

Этот результат полезен для матричных многообразий. Например, многообразие Stiefel
\[
\St(p,n)=\{X\in\R^{n\times p}:X^\top X=I_p\}
\]
задается системой гладких уравнений. Дифференцируя ограничение $X^\top X=I_p$, получаем касательное пространство
\[
T_X\St(p,n)=\{\Xi\in\R^{n\times p}:X^\top\Xi+\Xi^\top X=0\}.
\]
Многообразие Grassmann $\Gr(p,n)$ устроено немного иначе: его точка --- это не конкретный ортонормированный базис, а само $p$-мерное подпространство. Эти два примера важны для интуиции, потому что матрица фиксированного ранга фактически содержит два подпространства: столбцовое и строковое.

\begin{definition}
Риманова метрика на многообразии задает скалярное произведение на каждом касательном пространстве $T_xM$. В случае вложенного многообразия $M\subset \R^N$ чаще всего используется метрика, индуцированная обычным евклидовым скалярным произведением.
\end{definition}

Пусть $f:M\to\R$ --- функция, которую нужно минимизировать. Риманов градиент $\grad f(x)$ определяется условием
\[
df_x[\xi] = \langle \grad f(x), \xi \rangle_x
\qquad \text{для всех } \xi \in T_xM.
\]
Если $M$ вложено в евклидово пространство и используется индуцированная метрика, то риманов градиент получается особенно просто:
\begin{equation}
\label{eq:riemannian-gradient-projection}
\grad f(x) = \Proj_{T_xM}\bigl(\nabla \bar f(x)\bigr),
\end{equation}
где $\bar f$ --- продолжение функции на объемлющее пространство, а $\Proj_{T_xM}$ --- ортогональная проекция на касательное пространство.

После шага вдоль касательного направления точка обычно снова не лежит на многообразии. Поэтому нужен способ вернуться обратно.
\begin{definition}
Ретракцией называется отображение $\Retr_x:T_xM\to M$, которое переводит касательный вектор в точку многообразия и локально совпадает с геодезическим шагом с точностью первого порядка:
\[
\Retr_x(0_x)=x, \qquad D\Retr_x(0_x)=\mathrm{id}_{T_xM}.
\]
\end{definition}

Именно из этих объектов складывается общий риманов градиентный спуск:
\[
x_{k+1}=\Retr_{x_k}\bigl(-\eta_k \grad f(x_k)\bigr).
\]

\subsection{Базовые постановки задачи}

Для сравнения с геометрическим подходом удобно выписать несколько базовых моделей. Первая --- выпуклая регуляризованная постановка
\begin{equation}
\label{eq:convex}
\min_{X \in \R^{m\times n}}
\frac12 \fro{\Proj_\Omega(X-Y)}^2 + \lambda \nuc{X}.
\end{equation}

Здесь первое слагаемое отвечает за согласование с наблюдаемыми элементами, а второе поощряет низкий ранг. Эта модель удобна тем, что она выпукла, но ее вычислительная реализация в больших задачах может быть дорогой.

Вторая постановка --- задача на множестве матриц фиксированного ранга:
\begin{equation}
\label{eq:fixed-rank}
\min_{X \in \R^{m\times n}} \frac12 \fro{\Proj_\Omega(X-Y)}^2
\qquad \text{при условии} \qquad \rank(X)=r.
\end{equation}

Задача \eqref{eq:fixed-rank} уже невыпукла, но именно она является естественной для геометрических методов. Если ранг выбирается заранее, то в этой модели мы работаем напрямую с низкоранговой структурой, а не через ее выпуклую замену.

Для шумных данных полезно также рассмотреть регуляризованную fixed-rank постановку:
\begin{equation}
\label{eq:fixed-rank-regularized}
\min_{X \in \Mcal_r} f_\lambda(X)
= \frac12 \fro{\Proj_\Omega(X-Y)}^2 + \frac{\lambda}{2}\fro{X}^2,
\qquad \lambda\ge 0.
\end{equation}
Здесь штраф $\fro{X}^2$ не меняет само многообразие допустимых матриц, но меняет направление спуска. Интуитивно он не дает решению слишком сильно подстраиваться под шум в наблюдаемых элементах и особенно важен в случаях, когда выбранный ранг больше истинного.


\subsection{Многообразие матриц фиксированного ранга}

Рассмотрим множество
\[
\Mcal_r = \{X \in \R^{m \times n} : \rank(X)=r\}.
\]
Важно, что здесь используется условие ровно $\rank(X)=r$, а не $\rank(X)\le r$. Множество матриц ранга не выше $r$ содержит точки меньшего ранга, в которых меняется локальная размерность, поэтому оно уже имеет особенности. Для римановой оптимизации удобна именно гладкая страта фиксированного ранга.

\begin{proposition}
Множество $\Mcal_r$ является гладким вложенным многообразием в $\R^{m\times n}$ размерности
\[
\dim \Mcal_r=r(m+n-r).
\]
\end{proposition}

\begin{proof}
Любую матрицу $X\in\Mcal_r$ можно локально представить через факторизацию
\[
X=AB^\top,
\]
где $A\in\R^{m\times r}$ и $B\in\R^{n\times r}$ имеют полный столбцовый ранг. Такая параметризация имеет $mr+nr$ параметров, но она не единственна: для любой обратимой матрицы $Q\in\R^{r\times r}$ пары $(A,B)$ и $(AQ,BQ^{-\top})$ задают одну и ту же матрицу. Эта неоднозначность имеет $r^2$ степеней свободы. Поэтому локальная размерность равна
\[
mr+nr-r^2=r(m+n-r).
\]
Гладкость следует из того, что в окрестности матрицы ранга $r$ можно выбрать невырожденный минор размера $r\times r$ и использовать соответствующие блоки как локальные координаты.
\end{proof}

Если $X \in \Mcal_r$, то его удобно представлять в виде компактного SVD
\[
X = U\Sigma V^\top,
\]
где $U \in \R^{m\times r}$ и $V \in \R^{n\times r}$ имеют ортонормированные столбцы, а $\Sigma \in \R^{r\times r}$ невырождена. Такое представление особенно важно для вычислений: число параметров уменьшается с $mn$ до порядка $r(m+n)$.

\begin{proposition}
Если $X=U\Sigma V^\top\in\Mcal_r$, то касательное пространство к $\Mcal_r$ в точке $X$ можно записать как
\[
T_X\Mcal_r =
\left\{
UMV^\top + U_pV^\top + UV_p^\top :
M \in \R^{r\times r},\,
U^\top U_p = 0,\,
V^\top V_p = 0
\right\}.
\]
\end{proposition}

\begin{proof}
Касательное направление можно получить, дифференцируя представление $X(t)=U(t)\Sigma(t)V(t)^\top$ в точке $t=0$. Член $UMV^\top$ отвечает за изменение внутренней матрицы ранга $r$, а члены $U_pV^\top$ и $UV_p^\top$ отвечают за изменение столбцового и строкового подпространств. Условия $U^\top U_p=0$ и $V^\top V_p=0$ фиксируют именно движения, выходящие из текущих подпространств, и убирают лишнюю неоднозначность параметризации.
\end{proof}

\begin{lemma}
Пусть $P_U=UU^\top$ и $P_V=VV^\top$. Тогда ортогональная проекция произвольной матрицы $Z\in\R^{m\times n}$ на касательное пространство $T_X\Mcal_r$ имеет вид
\begin{equation}
\label{eq:tangent-projection-fixed-rank}
\Proj_{T_X\Mcal_r}(Z)=P_UZ+ZP_V-P_UZP_V.
\end{equation}
\end{lemma}

\begin{proof}
Разложим пространство матриц по столбцовому пространству $U$ и его ортогональному дополнению, а также по строковому пространству $V$ и его ортогональному дополнению. Тогда любая матрица $Z$ раскладывается на четыре блока:
\[
Z=P_UZP_V+P_UZ(I-P_V)+(I-P_U)ZP_V+(I-P_U)Z(I-P_V).
\]
Первые три слагаемых лежат в $T_X\Mcal_r$: они соответствуют изменению внутреннего блока, строкового и столбцового подпространств. Последнее слагаемое одновременно ортогонально столбцовому и строковому подпространствам, поэтому является нормальной компонентой. Удаляя эту нормальную компоненту, получаем
\[
Z-(I-P_U)Z(I-P_V)=P_UZ+ZP_V-P_UZP_V.
\]
Это и есть ортогональная проекция на касательное пространство.
\end{proof}

Эта формула важна не только теоретически: именно она используется в реализации риманова градиентного спуска.

\subsection{Евклидов и риманов градиенты}

Рассмотрим функцию потерь
\[
f(X) := \frac12 \fro{\Proj_\Omega(X-Y)}^2.
\]

\begin{lemma}
Евклидов градиент функции $f$ имеет вид
\[
\nabla f(X) = \Proj_\Omega(X-Y).
\]
\end{lemma}

\begin{proof}
Для любого приращения $H \in \R^{m\times n}$ имеем
\[
f(X+H)-f(X)
= \frac12 \fro{\Proj_\Omega(X+H-Y)}^2
 - \frac12 \fro{\Proj_\Omega(X-Y)}^2.
\]
Линейный по $H$ член в этом выражении равен
\[
\langle \Proj_\Omega(X-Y), \Proj_\Omega(H)\rangle.
\]
Поскольку $\Proj_\Omega$ является ортогональной проекцией, получаем
\[
\langle \Proj_\Omega(X-Y), H\rangle,
\]
откуда и следует формула для градиента.
\end{proof}

Для регуляризованной функции из \eqref{eq:fixed-rank-regularized} евклидов градиент меняется очень просто:
\[
\nabla f_\lambda(X)=\Proj_\Omega(X-Y)+\lambda X.
\]
Первое слагаемое отвечает за согласование с наблюдениями, а второе мягко уменьшает норму решения. Поэтому регуляризация почти не усложняет алгоритм: нужно только добавить $\lambda X$ к обычному евклидову градиенту.

На многообразии $\Mcal_r$ используется не полный евклидов градиент, а его проекция на касательное пространство:
\[
\grad_{\Mcal_r} f(X) = \Proj_{T_X\Mcal_r}\bigl(\Proj_\Omega(X-Y)\bigr).
\]
Для регуляризованной задачи соответственно получаем
\[
\grad_{\Mcal_r} f_\lambda(X)
= \Proj_{T_X\Mcal_r}\bigl(\Proj_\Omega(X-Y)+\lambda X\bigr).
\]
Таким образом, геометрия задачи входит в алгоритм уже на уровне самого направления спуска.

\subsection{Риманов градиентный спуск}

Общий шаг риманова градиентного спуска записывается в виде
\begin{equation}
\label{eq:rgd}
X_{k+1} = R_{X_k}\bigl(-\eta_k \grad_{\Mcal_r} f(X_k)\bigr),
\end{equation}
где $\eta_k > 0$ --- длина шага, а $R_{X_k}$ --- ретракция, возвращающая точку из касательного пространства обратно на многообразие ранга $r$.

В численной реализации удобно использовать ретракцию через усеченное SVD:
\begin{equation}
\label{eq:fixed-rank-retraction}
R_X(\Xi)=\Pi_r(X+\Xi),
\end{equation}
где $\Pi_r$ обозначает лучшую аппроксимацию ранга $r$ в норме Фробениуса.

\begin{proposition}
Отображение $R_X(\Xi)=\Pi_r(X+\Xi)$ является естественной ретракцией на многообразии $\Mcal_r$ в окрестности точки $X\in\Mcal_r$.
\end{proposition}

\begin{proof}
Во-первых, $R_X(0)=\Pi_r(X)=X$, так как сама матрица $X$ уже имеет ранг $r$. Во-вторых, по теореме Эккарта--Янга--Мирского усеченное SVD дает ближайшую к $X+\Xi$ матрицу ранга $r$ в норме Фробениуса. Поэтому после малого касательного шага $X+\Xi$, который вообще говоря уже не обязан иметь ранг ровно $r$, оператор $\Pi_r$ возвращает ближайшую точку на низкоранговое множество.

Остается понять условие первого порядка. Если $\Xi\in T_X\Mcal_r$, то по определению касательного пространства существует гладкая кривая $\gamma(t)\subset\Mcal_r$ такая, что $\gamma(0)=X$ и $\dot\gamma(0)=\Xi$. Значит,
\[
\gamma(t)=X+t\Xi+O(t^2).
\]
Матрица $X+t\Xi$ находится от многообразия $\Mcal_r$ на расстоянии порядка $O(t^2)$, потому что рядом есть точка $\gamma(t)$ ранга $r$. Так как $\Pi_r$ выбирает ближайшую матрицу ранга $r$, получаем
\[
\Pi_r(X+t\Xi)=X+t\Xi+O(t^2).
\]
Следовательно,
\[
DR_X(0)[\Xi]=\left.\frac{d}{dt}\right|_{t=0}\Pi_r(X+t\Xi)=\Xi.
\]
Это и есть условие $DR_X(0)=\mathrm{id}_{T_X\Mcal_r}$, поэтому отображение является ретракцией.
\end{proof}

Такая ретракция естественна для задачи matrix completion: она сохраняет фиксированный ранг после каждого шага и использует стандартную операцию усеченного SVD, которая уже лежит в основе многих низкоранговых методов \cite{vandereycken}.

С точки зрения реализации схема алгоритма выглядит так:
\begin{enumerate}[label=\arabic*.]
    \item выбрать ранг $r$ и начальное приближение $X_0 \in \Mcal_r$;
    \item вычислить евклидов градиент $\Proj_\Omega(X_k-Y)$;
    \item спроецировать его на касательное пространство $T_{X_k}\Mcal_r$;
    \item выбрать шаг $\eta_k$;
    \item выполнить ретракцию и получить новую точку $X_{k+1}$.
\end{enumerate}

На формальном уровне это очень похоже на обычный градиентный спуск. Однако принципиальное отличие состоит в том, что все обновления согласованы с геометрией множества $\Mcal_r$. Благодаря этому ограничение на ранг выполняется на каждой итерации, а сама структура задачи используется явно, а не косвенно.

Возникает естественный вопрос: почему вместо ретракции не идти по точной геодезической? В римановой геометрии такой шаг задавался бы экспоненциальным отображением. Однако для многообразия матриц фиксированного ранга точная геодезика вычисляется существенно сложнее, чем усеченное SVD или обновление низкоранговых факторов. Для градиентного спуска это обычно избыточно: стандартная теория римановой оптимизации требует не точного геодезического шага, а корректного возвращения на многообразие с совпадением первого порядка. Ретракция как раз удовлетворяет этому требованию, но при этом заметно проще и дешевле вычислительно. Поэтому в численных методах для matrix completion почти всегда используют ретракции, а не точные геодезики.

\subsection{Компактная реализация через $X=USV^\top$}

Ретракция \eqref{eq:fixed-rank-retraction} удобна для теории, но в прямой реализации она требует усеченного SVD матрицы размера $m\times n$ на каждой итерации. Для больших матриц это может быть именно тем узким местом, которого мы хотим избежать. Поэтому в практической реализации полезно хранить не всю матрицу $X$, а ее компактные факторы
\[
X=USV^\top,\qquad U^\top U=I_r,\quad V^\top V=I_r,
\]
где $U\in\R^{m\times r}$, $V\in\R^{n\times r}$, а $S\in\R^{r\times r}$ невырождена.

Важно подчеркнуть, что это не превращает метод в обычный покоординатный градиентный спуск по $U$, $S$ и $V$. Мы по-прежнему делаем шаг по матрице $X$ на многообразии $\Mcal_r$, но записываем этот шаг в низкоранговых координатах. Поэтому факторы $U,S,V$ используются как удобный способ хранить точку и выполнять ретракцию.

\subsubsection{Разложение риманова градиента}

Пусть
\[
G=\nabla f_\lambda(X)=\Proj_\Omega(X-Y)+\lambda X
\]
--- евклидов градиент регуляризованной функции. Проекцию $G$ на касательное пространство можно хранить через три компоненты:
\[
M=U^\top G V,
\]
\[
U_\perp=GV-UM,\qquad V_\perp=G^\top U-VM^\top.
\]
В этих обозначениях
\[
\grad_{\Mcal_r} f_\lambda(X)
=UMV^\top+U_\perp V^\top+UV_\perp^\top,
\]
причем $U^\top U_\perp=0$ и $V^\top V_\perp=0$. Это ровно та же проекция на касательное пространство, что и в формуле \eqref{eq:tangent-projection-fixed-rank}, только записанная в факторизованном виде.

Содержательно эти три компоненты отвечают за разные части движения. Матрица $M$ меняет внутреннее ядро $S$, компонент $U_\perp$ поворачивает столбцовое подпространство, а компонент $V_\perp$ поворачивает строковое подпространство.

\subsubsection{Касательный шаг}

Пусть длина шага равна $\eta>0$. Тогда шаг в направлении минус риманова градиента имеет вид
\[
\Delta M=-\eta M,\qquad
\Delta U=-\eta U_\perp,\qquad
\Delta V=-\eta V_\perp.
\]
После такого шага получаем матрицу
\begin{equation}
\label{eq:compact-tangent-step}
X+\Xi
=U(S+\Delta M)V^\top+\Delta U V^\top+U\Delta V^\top.
\end{equation}
Эта матрица уже не обязана иметь ранг ровно $r$. Поэтому, как и раньше, нужен возврат на многообразие, то есть ретракция.

\subsubsection{Ретракция через малое ядро}

В прямой реализации можно было бы взять матрицу \eqref{eq:compact-tangent-step} целиком и выполнить для нее усеченное SVD. Но делать этого не нужно. Заметим, что $\Delta U$ и $\Delta V$ имеют только $r$ столбцов, поэтому их можно разложить с помощью QR:
\[
\Delta U=Q_U R_U,\qquad \Delta V=Q_V R_V.
\]
Тогда касательный шаг раскладывается следующим образом:
\[
X+\Xi
=
\begin{bmatrix}U & Q_U\end{bmatrix}
\begin{bmatrix}
S+\Delta M & R_V^\top\\
R_U & 0
\end{bmatrix}
\begin{bmatrix}V & Q_V\end{bmatrix}^\top.
\]
Действительно, если раскрыть произведение, то получим
\[
U(S+\Delta M)V^\top+Q_UR_UV^\top+UR_V^\top Q_V^\top,
\]
а это ровно
\[
U(S+\Delta M)V^\top+\Delta U V^\top+U\Delta V^\top.
\]

Теперь большая матрица $X+\Xi$ представлена через малое ядро
\[
K=
\begin{bmatrix}
S+\Delta M & R_V^\top\\
R_U & 0
\end{bmatrix},
\]
размер которого не больше $2r\times 2r$. Поэтому для ретракции достаточно сделать SVD только этой маленькой матрицы:
\[
K=\widetilde U\widetilde S\widetilde V^\top.
\]
После этого оставляем первые $r$ сингулярных компонент и получаем новые факторы:
\[
U_+=\begin{bmatrix}U & Q_U\end{bmatrix}\widetilde U_r,\qquad
S_+=\widetilde S_r,\qquad
V_+=\begin{bmatrix}V & Q_V\end{bmatrix}\widetilde V_r.
\]
Итоговая новая точка равна
\[
X_+=U_+S_+V_+^\top.
\]

Эта процедура задает ту же идею ретракции: после касательного шага мы снова получаем матрицу ранга $r$. Но вычислительно она лучше согласована с целью работы: вместо полного SVD большой матрицы на каждой итерации используются QR-разложения тонких матриц размера $m\times r$ и $n\times r$ и SVD малого ядра размера не больше $2r\times 2r$.

Схема компактного шага выглядит так:
\begin{enumerate}[label=\arabic*.]
    \item представить текущую точку как $X_k=U_kS_kV_k^\top$;
    \item вычислить $G_k=\Proj_\Omega(X_k-Y)+\lambda X_k$;
    \item найти компоненты $M_k$, $U_{\perp,k}$, $V_{\perp,k}$ риманова градиента;
    \item выбрать шаг $\eta_k$ с помощью backtracking line search;
    \item построить $\Delta M_k=-\eta_kM_k$, $\Delta U_k=-\eta_kU_{\perp,k}$, $\Delta V_k=-\eta_kV_{\perp,k}$;
    \item выполнить QR-разложения $\Delta U_k=Q_UR_U$ и $\Delta V_k=Q_VR_V$;
    \item собрать малое ядро $K_k$ и сделать его SVD;
    \item получить новые факторы $U_{k+1}$, $S_{k+1}$, $V_{k+1}$.
\end{enumerate}

\begin{algorithm}[H]
\caption{Компактный риманов градиентный спуск на $\Mcal_r$}
\label{alg:compact-rgd}
\begin{algorithmic}[1]
\Require наблюдаемая матрица $Y$, маска $\Omega$, ранг $r$, параметр $\lambda$, начальные факторы $U_0,S_0,V_0$
\Ensure приближение $\widehat X=USV^\top$ ранга $r$
\For{$k=0,1,2,\ldots$}
    \State $X_k \gets U_kS_kV_k^\top$
    \State $G_k \gets \Proj_\Omega(X_k-Y)+\lambda X_k$
    \State $M_k \gets U_k^\top G_kV_k$
    \State $U_{\perp,k} \gets G_kV_k-U_kM_k$
    \State $V_{\perp,k} \gets G_k^\top U_k-V_kM_k^\top$
    \State выбрать $\eta_k$ с помощью backtracking line search
    \State $\Delta M_k\gets-\eta_kM_k$, $\Delta U_k\gets-\eta_kU_{\perp,k}$, $\Delta V_k\gets-\eta_kV_{\perp,k}$
    \State выполнить QR-разложения $\Delta U_k=Q_UR_U$, $\Delta V_k=Q_VR_V$
    \State $K_k\gets
    \begin{bmatrix}
    S_k+\Delta M_k & R_V^\top\\
    R_U & 0
    \end{bmatrix}$
    \State вычислить SVD малого ядра $K_k=\widetilde U\widetilde S\widetilde V^\top$
    \State $U_{k+1}\gets [U_k,Q_U]\widetilde U_r$
    \State $S_{k+1}\gets \widetilde S_r$
    \State $V_{k+1}\gets [V_k,Q_V]\widetilde V_r$
    \If{выполнен критерий остановки}
        \State \textbf{break}
    \EndIf
\EndFor
\State \Return $\widehat X=U_kS_kV_k^\top$
\end{algorithmic}
\end{algorithm}

Поэтому в практической части мы оставляем прямую SVD-ретракцию как прозрачный прототип, но дополнительно реализуем компактный вариант $USV^\top$ как основной кандидат для масштабирования.

\subsection{Почему этот подход интересен для нашей работы}

Сравнение выпуклой модели \eqref{eq:convex} и фиксированно-ранговой модели \eqref{eq:fixed-rank} показывает, что между ними есть не только техническое, но и содержательное различие.

\begin{table}[H]
\centering
\caption{Сравнение двух классов методов}
\label{tab:compare}
\begin{tabular}{p{0.30\linewidth} p{0.28\linewidth} p{0.28\linewidth}}
\toprule
\textbf{Критерий} & \textbf{Ядерная норма} & \textbf{Фиксированный ранг / риманов подход} \\
\midrule
Теоретическая прозрачность & Высокая, сильные гарантии & Ниже в базовой форме, больше зависит от алгоритма \\
\addlinespace
Выпуклость & Да & Нет \\
\addlinespace
Работа с рангом & Через релаксацию и регуляризацию & Напрямую через ограничение $\rank(X)=r$ \\
\addlinespace
Масштабируемость & Может быть ограничена SVD & Часто лучше при малом ранге \\
\addlinespace
Чувствительность к старту & Обычно ниже & Может быть существенной \\
\addlinespace
Исследовательский интерес & Гарантии восстановления & Поведение алгоритма, устойчивость, сходимость \\
\bottomrule
\end{tabular}
\end{table}

Для нашей работы особенно важно последнее различие. Если выпуклые методы обычно сравниваются с точки зрения теории и итоговой ошибки, то для геометрических алгоритмов возникают дополнительные вопросы: как влияет стартовая точка, как быстро метод сходится, насколько устойчив он к шуму и в каких режимах оказывается лучше конкурентов. Именно эти вопросы и превращают тему курсовой работы в потенциальный мини-ресерч.

\fi

\newpage


\section{Практическая часть}
\label{sec:empirical}

\subsection{Код и постановка эксперимента}

Для практической части написан отдельный Python-скрипт:
\begin{center}
\path{matrix_completion_experiments.py}
\end{center}
Скрипт генерирует низкоранговые матрицы, скрывает часть элементов, добавляет шум и сравнивает методы восстановления. Такой синтетический эксперимент нужен как контролируемая проверка: мы знаем настоящую матрицу $X_*$ и можем честно измерять ошибку.

Матрица строится по модели
\[
X_* = AB^\top,
\]
где $A\in\R^{m\times r}$ и $B\in\R^{n\times r}$. Затем часть элементов отдается алгоритмам для обучения, а часть известных элементов откладывается как test. На наблюдаемых позициях берется
\[
Y_{ij}=(X_*)_{ij}+\varepsilon_{ij}, \qquad (i,j)\in\Omega,
\]
где $\varepsilon_{ij}$ отвечает за шум.

\subsection{Методы и метрики}

В коде сравниваются ALS, Soft-Impute, риманов градиентный спуск с прямой SVD-ретракцией и компактный риманов спуск в представлении $X=USV^\top$. Для риманова метода также добавлена L2-регуляризация из \eqref{eq:fixed-rank-regularized}. Она нужна потому, что при шуме и завышенном ранге обычный fixed-rank метод может слишком сильно подстраиваться под наблюдаемые элементы.

Основные метрики:
\begin{enumerate}[label=\arabic*.]
    \item RMSE и MAE на test-элементах;
    \item относительная ошибка $\fro{\widehat X-X_*}/\fro{X_*}$;
    \item время работы;
    \item число итераций.
\end{enumerate}

Для параметра регуляризации $\lambda$ используются два режима. Первый --- перебор сетки, чтобы увидеть, какие значения вообще работают лучше. Второй --- validation-подбор: часть наблюдаемых элементов временно откладывается, $\lambda$ выбирается по ним, а итоговая модель затем заново обучается на всех наблюдаемых элементах.

\begin{algorithm}[H]
\caption{Validation-подбор параметра $\lambda$}
\label{alg:validation-lambda}
\begin{algorithmic}[1]
\Require наблюдаемые индексы $\Omega_{\mathrm{train}}$, сетка $\Lambda$, доля validation $\rho$
\Ensure выбранное значение $\lambda_*$ и итоговая матрица $\widehat X$
\State разбить $\Omega_{\mathrm{train}}$ на $\Omega_{\mathrm{fit}}$ и $\Omega_{\mathrm{val}}$
\For{$\lambda\in\Lambda$}
    \State обучить RGD на $\Omega_{\mathrm{fit}}$ с параметром $\lambda$
    \State посчитать $\operatorname{RMSE}_{\mathrm{val}}(\lambda)$ на $\Omega_{\mathrm{val}}$
\EndFor
\State $\lambda_*\gets \arg\min_{\lambda\in\Lambda}\operatorname{RMSE}_{\mathrm{val}}(\lambda)$
\State заново обучить RGD с $\lambda_*$ на всех индексах $\Omega_{\mathrm{train}}$
\State оценить итоговую ошибку на test-множестве
\State \Return $\lambda_*$ и $\widehat X$
\end{algorithmic}
\end{algorithm}

\subsection{Что меняется в экспериментах}

В сериях экспериментов меняются:
\begin{enumerate}[label=\arabic*.]
    \item доля наблюдаемых элементов;
    \item уровень шума;
    \item тип инициализации;
    \item правильный, заниженный или завышенный ранг;
    \item значение L2-регуляризации;
    \item прямая SVD-ретракция против компактного шага $USV^\top$.
\end{enumerate}

\subsection{Результаты на синтетических данных}

В текущем запуске использовались матрицы размера $60\times 60$ истинного ранга $4$. Каждый сценарий запускался в трех повторениях. В таблицах ниже оставлены методы, которые удобно сравнивать содержательно: ALS, Soft-Impute, компактный RGD без регуляризации, компактный RGD с несколькими значениями $\lambda$ и компактный RGD с validation-подбором $\lambda$.

Логика получилась последовательной. Сначала была реализована простая схема риманова градиентного спуска с SVD-ретракцией. Она хорошо работала при малом шуме, но хуже переносила шум, чем Soft-Impute. Потом была добавлена L2-регуляризация, после этого --- сетка по $\lambda$, затем validation-подбор, а в конце --- компактная реализация $USV^\top$, которая уменьшает стоимость шага.

\input{results_compact_l2_validation_r3/course_tables/course_research_progression.tex}

\input{results_compact_l2_validation_r3/course_tables/course_main_scalable_comparison.tex}

По таблице \ref{tab:main-scalable-results} видно, что компактный RGD с $\lambda=10^{-2}$ и $\lambda=3\cdot 10^{-2}$ дает меньшую среднюю test-ошибку, чем Soft-Impute, и при этом работает быстрее. Validation-подбор дает еще меньшую ошибку, но стоит дороже, потому что для каждого сценария приходится обучать несколько моделей.

Прямую SVD-ретракцию полезно оставить как ориентир качества. Это не основной вариант для масштабирования, но он показывает, насколько далеко можно продвинуться при более дорогом шаге.

\input{results_compact_l2_validation_r3/course_tables/course_quality_upper_bound.tex}

Таблица \ref{tab:validation-selected-upper-bound} показывает, что компактная реализация почти не теряет качество относительно прямой SVD-ретракции, но работает существенно быстрее. Поэтому дальше основной практический интерес представляет именно вариант $USV^\top$.

\input{results_compact_l2_validation_r3/course_tables/course_noise_scalable_comparison.tex}

Таблица \ref{tab:noise-scalable-results} показывает главный слабый участок метода. При нулевом и умеренном шуме компактный RGD выигрывает у Soft-Impute. При сильном шуме Soft-Impute становится устойчивее или дает очень близкий результат. Это ожидаемо: ядерная норма сильнее сглаживает сингулярные значения.

\input{results_compact_l2_validation_r3/course_tables/course_rank_scalable_comparison.tex}

По таблице \ref{tab:rank-scalable-results} видно, что RGD особенно хорошо работает, когда ранг выбран правильно или немного занижен. При завышенном ранге регуляризация становится важнее, иначе метод начинает подстраиваться под шум.

\input{results_compact_l2_validation_r3/course_tables/course_compact_lambda_sweep.tex}

Таблица \ref{tab:compact-lambda-sweep} показывает, что слишком маленькие $\lambda$ почти не помогают, а слишком большое $\lambda=0.1$ уже ухудшает качество. В текущей сетке лучше всего выглядят значения около $10^{-2}$ и $3\cdot 10^{-2}$.

\input{results_compact_l2_validation_r3/course_tables/course_compact_selected_lambda_counts.tex}

Validation чаще выбирает $\lambda$ порядка $10^{-2}$, но иногда предпочитает более сильную регуляризацию. Поэтому улучшение качества не является подглядыванием в test, а получается через отдельную процедуру выбора гиперпараметра.

\input{results_compact_l2_validation_r3/course_tables/course_compact_validation_wins.tex}

Итог по синтетике такой: риманов градиентный спуск хорошо выглядит, когда данные действительно близки к низкоранговым и шум не слишком большой. Компактная реализация делает метод быстрее, а L2-регуляризация заметно улучшает устойчивость.

\subsection{Следующий шаг}

После синтетических экспериментов нужно перейти к экономическим или панельным данным. Часть известных значений будет скрываться искусственно, чтобы методы можно было сравнить по восстановлению этих значений. При этом реальные пропуски могут быть неслучайными, поэтому прикладной эксперимент лучше рассматривать как дополнительную проверку, а не как замену синтетике.

\newpage


\section*{Заключение}
\phantomsection
\addcontentsline{toc}{section}{Заключение}

\textbf{Выводы.} В предварительном тексте сформирована основа курсовой работы по геометрическим методам для заполнения матриц. Во введении уточнена исследовательская постановка, в обзоре литературы показан переход от выпуклых методов к fixed-rank и римановым подходам, а в теоретической части добавлены необходимые элементы дифференциальной геометрии: касательное пространство, риманов градиент, проекция и ретракция.

\textbf{Основной результат текущего этапа.} Работа уже имеет не только теоретический, но и практический каркас. Подготовлен код для синтетических экспериментов, где можно отдельно менять шум, долю наблюдений, инициализацию, выбранный ранг и уровень L2-регуляризации в римановом методе. Также добавлен validation-подбор $\lambda$, который отделяет исследовательский oracle-анализ от честной процедуры выбора гиперпараметра. Отдельно реализована компактная версия риманова спуска в факторах $USV^\top$, что позволяет проверить не только качество восстановления, но и вычислительный выигрыш от работы в низкоранговых координатах. Это позволяет дальше переходить от описания метода к полноценному сравнению алгоритмов.

\textbf{Дальнейшие шаги.} На следующем этапе необходимо провести основную серию экспериментов по сетке значений $\lambda$, сравнить прямую и компактную реализации риманова метода с Soft-Impute, подготовить графики сходимости и перейти к проверке на прикладных экономических данных. После этого можно будет формулировать уже не только методические, но и исследовательские выводы о режимах, в которых риманов градиентный спуск оказывается наиболее полезным.

\begin{thebibliography}{99}
    \phantomsection
    \addcontentsline{toc}{section}{Список литературы}

\bibitem{ivanov-dg}
А. О. Иванов.
\newblock Классическая дифференциальная геометрия; Дифференциальная геометрия и топология.
\newblock Лекционные материалы курса, мехмат МГУ / Teach-in.
\newblock URL: \url{https://teach-in.ru/course/differential-geometry-and-tensor-analysis-ivanov};
\url{https://teach-in.ru/course/differential-geometry-and-tensor-analysis-ivanov-p2}.

\bibitem{fomenko-dg}
А. Т. Фоменко.
\newblock Классическая дифференциальная геометрия.
\newblock Лекционные материалы курса, мехмат МГУ / Teach-in.
\newblock URL: \url{https://teach-in.ru/course/classical-difgeom-fomenko/about}.

\bibitem{penskoy-seminars}
А. В. Пенской.
\newblock Классическая дифференциальная геометрия: семинары.
\newblock Семинарские материалы курса, мехмат МГУ / Teach-in.
\newblock URL: \url{https://teach-in.ru/course/kdg-seminars-penskoy/about}.

\bibitem{candes-recht}
Emmanuel J. Cand\`es, Benjamin Recht.
\newblock Exact Matrix Completion via Convex Optimization.
\newblock \textit{arXiv preprint} arXiv:0805.4471, 2008.
\newblock URL: \url{https://arxiv.org/abs/0805.4471}.

\bibitem{cai-svt}
Jian-Feng Cai, Emmanuel J. Cand\`es, Zuowei Shen.
\newblock A Singular Value Thresholding Algorithm for Matrix Completion.
\newblock \textit{arXiv preprint} arXiv:0810.3286, 2008.
\newblock URL: \url{https://arxiv.org/abs/0810.3286}.

\bibitem{candes-tao}
Emmanuel J. Cand\`es, Terence Tao.
\newblock The Power of Convex Relaxation: Near-Optimal Matrix Completion.
\newblock \textit{arXiv preprint} arXiv:0903.1476, 2009.
\newblock URL: \url{https://arxiv.org/abs/0903.1476}.

\bibitem{recht}
Benjamin Recht.
\newblock A Simpler Approach to Matrix Completion.
\newblock \textit{Journal of Machine Learning Research}, 12:3413--3430, 2011.
\newblock URL: \url{https://www.jmlr.org/papers/v12/recht11a.html}.

\bibitem{keshavan}
Raghunandan H. Keshavan, Andrea Montanari, Sewoong Oh.
\newblock Matrix Completion from a Few Entries.
\newblock \textit{arXiv preprint} arXiv:0901.3150, 2009.
\newblock URL: \url{https://arxiv.org/abs/0901.3150}.

\bibitem{candes-plan}
Emmanuel J. Cand\`es, Yaniv Plan.
\newblock Matrix Completion with Noise.
\newblock \textit{arXiv preprint} arXiv:0903.3131, 2009.
\newblock URL: \url{https://arxiv.org/abs/0903.3131}.

\bibitem{mazumder}
Rahul Mazumder, Trevor Hastie, Robert Tibshirani.
\newblock Spectral Regularization Algorithms for Learning Large Incomplete Matrices.
\newblock \textit{Journal of Machine Learning Research}, 11:2287--2322, 2010.
\newblock URL: \url{https://www.jmlr.org/papers/v11/mazumder10a.html}.

\bibitem{inductive}
Prateek Jain, Inderjit S. Dhillon.
\newblock Provable Inductive Matrix Completion.
\newblock \textit{arXiv preprint} arXiv:1306.0626, 2013.
\newblock URL: \url{https://arxiv.org/abs/1306.0626}.

\bibitem{overview}
Mark A. Davenport, Justin Romberg.
\newblock An Overview of Low-Rank Matrix Recovery from Incomplete Observations.
\newblock \textit{arXiv preprint} arXiv:1601.06422, 2016.
\newblock URL: \url{https://arxiv.org/abs/1601.06422}.

\bibitem{vandereycken}
Bart Vandereycken.
\newblock Low-rank Matrix Completion by Riemannian Optimization --- extended version.
\newblock \textit{arXiv preprint} arXiv:1209.3834, 2012.
\newblock URL: \url{https://arxiv.org/abs/1209.3834}.

\bibitem{r3mc}
B. Mishra, R. Sepulchre.
\newblock R3MC: A Riemannian Three-Factor Algorithm for Low-Rank Matrix Completion.
\newblock \textit{arXiv preprint} arXiv:1306.2672, 2013.
\newblock URL: \url{https://arxiv.org/abs/1306.2672}.

\bibitem{causal-panel}
Susan Athey, Mohsen Bayati, Nikolay Doudchenko, Guido W. Imbens, Khashayar Khosravi.
\newblock Matrix Completion Methods for Causal Panel Data Models.
\newblock \textit{arXiv preprint} arXiv:1710.10251, 2017.
\newblock URL: \url{https://arxiv.org/abs/1710.10251}.

\end{thebibliography}

\end{document}
